% Standalone resolution manuscript generated from solution.md.
% Compile with XeLaTeX; author and review metadata are maintained in Markdown.
\documentclass[10pt,a4paper]{article}
\usepackage[margin=24mm,headheight=15pt]{geometry}
\usepackage{fontspec}
\setmainfont{DejaVuSerif.ttf}[BoldFont=DejaVuSerif-Bold.ttf,ItalicFont=DejaVuSerif-Italic.ttf,BoldItalicFont=DejaVuSerif-BoldItalic.ttf]
\setsansfont{DejaVuSans.ttf}[BoldFont=DejaVuSans-Bold.ttf,ItalicFont=DejaVuSans-Oblique.ttf]
\setmonofont{DejaVuSansMono.ttf}
\usepackage{amsmath,amssymb,unicode-math}
\setmathfont{latinmodern-math.otf}
\usepackage{xcolor,microtype,parskip,needspace,fancyhdr}
\usepackage{bookmark}
\usepackage{xurl}
\hypersetup{colorlinks=true,urlcolor=blue,linkcolor=blue,pdftitle={MF-02:
A uniform constant-factor bound for cubic composition
overhead},pdfauthor={George Stepaniants},pdflang={en-GB}}
\urlstyle{same}
\setlength{\emergencystretch}{3em}
\providecommand{\tightlist}{\setlength{\itemsep}{0pt}\setlength{\parskip}{0pt}}
\providecommand{\pandocbounded}[1]{#1}
\setcounter{secnumdepth}{0}
\allowdisplaybreaks[1]
\widowpenalty=10000
\clubpenalty=10000
\pagestyle{fancy}
\fancyhf{}
\fancyhead[L]{\small\sffamily RESEARCH MANUSCRIPT / MF-02}
\fancyhead[R]{\small\sffamily 12 September 2026}
\fancyfoot[L]{\parbox[b]{0.9\textwidth}{\fontsize{8}{10}\selectfont AI-assisted
manuscript; exact review and source-attribution evidence is retained in
the submission record.}}
\fancyfoot[R]{\footnotesize\thepage}
\begin{document}
{\raggedright\Large\bfseries MF-02: A uniform constant-factor bound for
cubic composition overhead\par}
\vspace{0.4em}
{\normalsize\bfseries George Stepaniants\par}
{\small Department of Computing and Mathematical Sciences, California
Institute of Technology, Pasadena, California, USA.\par}
\vspace{0.6em}
\pagestyle{plain}
\setlength{\abovedisplayskip}{6pt}
\setlength{\belowdisplayskip}{6pt}
\setlength{\abovedisplayshortskip}{4pt}
\setlength{\belowdisplayshortskip}{4pt}

This manuscript establishes a uniform constant-factor asymptotic bound
for the stage count in MF-02. The classical cubic iteration and the
prior constant-factor optimality results are credited in Section 6. The
explicit estimate proved here is at most one cubic stage per allowed
unrestricted multiplication. The exact smallest stage count and the
stronger same-budget optimization question remain open.

Prepared with substantial ChatGPT/Codex assistance at the author's
request. A separate
\href{https://github.com/ajt60gaibb/OpenProblemsInNLA/blob/main/references/stepaniants-mf02-2026-09-12/verification/independent-review/MF-02-independent-review.md}{independent
mathematical review} returned PASS for the theorem and the canonical
uniform asymptotic-order scope. A second
\href{https://github.com/ajt60gaibb/OpenProblemsInNLA/blob/main/references/stepaniants-mf02-2026-09-12/verification/source-review/REVIEW.md}{independent
source and scope review} confirms the attribution and the distinction
between the canonical asymptotic target and the stronger source
questions. This is an informal automated-agent audit, not external human
peer review or formal verification.

\subsection{1. Exact target and result}\label{exact-target-and-result}

Let \(I_\delta=[-1,-\delta]\cup[\delta,1]\), where \(0<\delta<1\). Let
\(\mathcal P_m\) be the real polynomials computed from \(1,x\), using
arbitrary real linear combinations and at most \(m\) nonscalar
multiplications. Define \[
E_m(\delta)=\inf_{p\in\mathcal P_m}
 \max_{x\in I_\delta}|p(x)-\operatorname{sign}(x)|.
\] For \(T\ge0\), let \[
C_T(\delta)=\inf_{a_t,b_t\in\mathbb R}
 \max_{x\in I_\delta}
 |(q_T\circ\cdots\circ q_1)(x)-\operatorname{sign}(x)|,
\qquad q_t(x)=a_tx+b_tx^3.
\] At \(T=0\) the composition is \(x\). The canonical MF-02 target asks
for the asymptotic dependence on both \(m,\delta\) of \[
T_{\min}(m,\delta)=
\inf\{T\in\mathbb N_0:C_T(\delta)\le E_m(\delta)\}.
\]

\textbf{Theorem.} For every \(0<\delta<1\), \[
T_{\min}(0,\delta)=T_{\min}(1,\delta)=1.
\] For every integer \(m\ge2\), \[
\left\lfloor\frac m2\right\rfloor
\le T_{\min}(m,\delta)\le m.
\tag{1}
\] Consequently, \[
\frac{m+1}{4}\le T_{\min}(m,\delta)\le m+1
\qquad(m\in\mathbb N_0,\ 0<\delta<1).
\tag{2}
\] Thus the asymptotic order is \(\Theta(m+1)\), with absolute constants
uniform over the entire gap range, including gaps depending on \(m\).

This theorem gives a constant-factor asymptotic comparison. It does not
give an exact formula for \(T_{\min}\), an optimal leading constant, or
the value of \(E_m\). The uniform two-sided order in (2) determines the
canonical asymptotic dependence on both parameters. The broader source
question about the exact smallest \(T\), and the comparison at the same
multiplication budget, remain unanswered by (1).

\Needspace{12\baselineskip}

\subsection{2. A degree-only lower
bound}\label{a-degree-only-lower-bound}

Put \[
r_\delta=\frac{1-\delta}{1+\delta}\in(0,1).
\]

\textbf{Lemma 1.} If \(p\in\mathbb R[x]\) has degree at most \(D\),
where \(D\ge1\) is an integer, then \[
\max_{x\in I_\delta}|p(x)-\operatorname{sign}(x)|
\ge r_\delta^D.
\tag{3}
\]

We first recall and prove the needed Chebyshev comparison. Let \(T_D\)
be the degree-\(D\) Chebyshev polynomial, characterized by \[
T_D(\cos\theta)=\cos(D\theta).
\] If \(R\) is a real polynomial of degree at most \(D\) with
\(|R|\le1\) on \([-1,1]\), then \[
|R(t)|\le |T_D(t)|\qquad(|t|>1).
\tag{4}
\] For \(t_0>1\), suppose otherwise. Replacing \(R\) by \(-R\) if
necessary, set \(c=R(t_0)/T_D(t_0)>1\). At the \(D+1\) points
\(\cos(j\pi/D)\), the polynomial \(R-cT_D\) has strict alternating
signs, because \(|R|\le1<c\). It has at least \(D\) distinct roots in
\((-1,1)\), as well as the root \(t_0\), contradicting its degree. The
polynomial is not identically zero because \(c>1\). The case \(t_0<-1\)
follows by replacing \(R(t)\) with \(R(-t)\) and using
\(T_D(-t)=(-1)^DT_D(t)\).

To prove (3), let the error of \(p\) be \(e\). If \(e\ge1\), the
assertion follows. The value \(e=0\) is impossible: agreement with \(1\)
throughout the interval \([\delta,1]\) would force \(p\equiv1\),
contradicting agreement with \(-1\) on the other interval. Hence
consider \(0<e<1\). The odd part \[
o(x)=\frac{p(x)-p(-x)}2
\] also has error at most \(e\), so \(1-e\le o(x)\le1+e\) on
\([\delta,1]\). As \(o\) is odd, \(o(x)^2=B(x^2)\) for a real polynomial
\(B\) of degree at most \(D\), and \(B(0)=0\). The polynomial \[
H(t)=\frac{1+e^2-B(t)}{2e}
\] has absolute value at most \(1\) on \([\delta^2,1]\). Map this
interval affinely to \([-1,1]\); the image of \(t=0\) is \[
-\frac{1+\delta^2}{1-\delta^2}.
\] By (4), \[
\frac{1+e^2}{2e}=H(0)
\le T_D\!\left(\frac{1+\delta^2}{1-\delta^2}\right)
=\frac{r_\delta^{-D}+r_\delta^D}{2}.
\tag{5}
\] The last equality follows from \(T_D(\cosh s)=\cosh(Ds)\) and \[
\frac{1+\delta^2}{1-\delta^2}
=\cosh\!\left(\log\frac{1+\delta}{1-\delta}\right).
\] The function \(e\mapsto(e+e^{-1})/2\) is strictly decreasing on
\((0,1)\). Thus (5) gives \(e\ge r_\delta^D\), as claimed.

Every polynomial in \(\mathcal P_m\) has degree at most \(2^m\): the
largest available degree starts at \(1\), linear combinations never
increase it, and each multiplication at most doubles it. This remains
true with stored-value reuse and an arbitrary number of free linear
combinations. Consequently, \[
E_m(\delta)\ge r_\delta^{\,2^m}>0.
\tag{6}
\] Likewise each length-\(T\) cubic composition has degree at most
\(3^T\), including degenerate stages, so \[
C_T(\delta)\ge r_\delta^{\,3^T}>0.
\tag{7}
\] Taking infima in these pointwise inequalities does not require
attainment.

\Needspace{12\baselineskip}

\subsection{3. One cubic step squares the interval error
ratio}\label{one-cubic-step-squares-the-interval-error-ratio}

For \(0<a<1\), set \[
A=1+a+a^2,\qquad
M=\frac{2A^{3/2}}{3\sqrt3},\qquad
g_a(x)=\frac{x(A-x^2)}M,
\] and define \[
\phi(a)=\frac{3\sqrt3\,a(1+a)}
                  {2(1+a+a^2)^{3/2}}.
\tag{8}
\] The polynomial \(g_a\) has the allowed form \(ux+vx^3\). It maps
\([a,1]\) into \([\phi(a),1]\). Indeed its unique critical point in that
interval is \(\sqrt{A/3}\): \[
A-3a^2=(1-a)(1+2a)>0,\qquad
3-A=(1-a)(a+2)>0.
\] At this point \(x(A-x^2)=M\), while the two endpoint values are both
\(a(1+a)\). The derivative is positive before the critical point and
negative after it. In particular \(0<\phi(a)<1\).

\textbf{Lemma 2.} For every \(0<a<1\), \[
\phi(a)\ge\frac{2a}{1+a^2},\qquad
\frac{1-\phi(a)}{1+\phi(a)}
\le \left(\frac{1-a}{1+a}\right)^2.
\tag{9}
\]

All quantities in the first comparison are positive. After cancelling
\(a\) and squaring, it is equivalent to \[
27(1+a)^2(1+a^2)^2-16(1+a+a^2)^3\ge0.
\] The exact identity \[
27(1+a)^2(1+a^2)^2-16(1+a+a^2)^3
=(1-a)^2(11a^4+28a^3+30a^2+28a+11)
\tag{10}
\] proves it. Since \(v\mapsto(1-v)/(1+v)\) is decreasing for \(v>0\),
the first comparison gives the second, using \[
\frac{1-2a/(1+a^2)}{1+2a/(1+a^2)}
=\left(\frac{1-a}{1+a}\right)^2.
\]

Starting with \(a_0=\delta\), let \(a_{t+1}=\phi(a_t)\). For \(T\ge1\),
the composition \[
G_T=g_{a_{T-1}}\circ\cdots\circ g_{a_0}
\] maps \([\delta,1]\) into \([a_T,1]\) and is odd. Scale its output by
\(2/(1+a_T)\), absorbing that scalar into the two coefficients of the
last cubic stage. The resulting allowed length-\(T\) composition has
error at most \[
\frac{1-a_T}{1+a_T}\le r_\delta^{\,2^T}.
\] Thus \[
C_T(\delta)\le r_\delta^{\,2^T}\qquad(T\ge1).
\tag{11}
\] No additional composition stage is used for the final scaling.
Coefficients may depend on \(T,\delta\), as permitted by the target.
Each stage requires at most two nonscalar multiplications.

\Needspace{12\baselineskip}

\subsection{4. Strict improvement survives taking
infima}\label{strict-improvement-survives-taking-infima}

\textbf{Lemma 3.} For fixed \(\delta\in(0,1)\), the sequence
\(C_T(\delta)\) is strictly decreasing in \(T\).

First \(C_0(\delta)=1-\delta\), and (11) implies \[
C_1(\delta)\le r_\delta^2<1-\delta=C_0(\delta).
\tag{12}
\] For \(T\ge1\), (7) and (11) give \(0<C_T(\delta)<1\). Take any
allowed length-\(T\) composition \(P\) with error \(e<1\). It is odd and
satisfies \(P([\delta,1])\subset[1-e,1+e]\). Set \(a=(1-e)/(1+e)\), so
that \[
\frac{P([\delta,1])}{1+e}\subset[a,1].
\] The polynomial in the single variable \(z\), \[
h_e(z)=\frac{2}{1+\phi(a)}
       g_a\!\left(\frac{z}{1+e}\right),
\] is one allowed cubic stage. Its composition with \(P\) has error at
most \[
\frac{1-\phi(a)}{1+\phi(a)}
\le\left(\frac{1-a}{1+a}\right)^2=e^2.
\] Choose a sequence of length-\(T\) compositions whose errors tend to
\(C_T(\delta)\); the errors are eventually less than \(1\). It follows
that \[
C_{T+1}(\delta)\le C_T(\delta)^2<C_T(\delta).
\tag{13}
\] This argument handles potentially unattained infima directly.

\Needspace{12\baselineskip}

\subsection{5. Proof of the theorem and
interpretation}\label{proof-of-the-theorem-and-interpretation}

For \(m\ge1\), combine (6) and (11) with \(T=m\): \[
C_m(\delta)\le r_\delta^{\,2^m}\le E_m(\delta).
\tag{14}
\] Hence the admissible set in the definition of \(T_{\min}\) is
nonempty and \(T_{\min}(m,\delta)\le m\).

For \(m\ge2\), put \(k=\lfloor m/2\rfloor\ge1\). Every length-\(k\)
cubic composition is computed using at most \(2k\le m\) multiplications,
so \[
E_m(\delta)\le C_k(\delta).
\] If \(T<k\), Lemma 3 yields \[
C_T(\delta)>C_k(\delta)\ge E_m(\delta).
\] Thus \(T_{\min}(m,\delta)\ge k\), establishing (1).

At \(m=0\), Lemma 1 with \(D=1\) and the linear polynomial
\(2x/(1+\delta)\) show \[
E_0(\delta)=r_\delta.
\] Every polynomial computable with one multiplication has degree at
most two. Its odd part is linear, so the same lower bound proves
\(E_1(\delta)=r_\delta\). Since \(C_0(\delta)=1-\delta>r_\delta\) and
\(C_1(\delta)\le r_\delta^2<r_\delta\), the asserted two endpoint values
follow. Finally, for \(m\ge2\), \(\lfloor m/2\rfloor\ge(m+1)/4\), and
the two small cases satisfy (2) as well.

The comparison uses at most \(2m\) matrix products for \(m\ge1\),
because an allowed cubic stage costs at most two. This is an upper
guarantee on multiplication overhead, not an assertion that factor two
is optimal or always necessary. The estimate remains uniform when
\(\delta\) tends to either endpoint at any rate with \(m\).

\Needspace{12\baselineskip}

\subsection{6. Attribution, scope, and
verification}\label{attribution-scope-and-verification}

The scaled cubic and its endpoint recurrence are due to Chen and Chow
(2014), Section 3, equations (3.3)-(3.6), printed pages 9-10; their
Theorems 3.1-3.2, printed page 11, analyze its convergence. Polar
Express, version 5, Appendix F, equation (17), gives the same cubic and
final centering. Its Theorem 3.1, equations (8)-(10), establishes
optimality within the fixed-degree composition class, with the Chen-Chow
antecedent credited in Section 3.2 and Appendix B. The construction in
Section 3 is not a new iteration.

Cheon, Kim and Kim (2020) already establish constant-factor
asymptotically optimal fixed-degree composition complexity, including
the cubic case. Their relevant locators are Section 2.1, printed page 9;
Lemma 3, page 14; Theorem 3, page 15; Table 1, page 17; and Section 3.4,
pages 17-18. A short synthesis of their global Lemma 3 and the classical
Chebyshev degree bound already gives the uniform constant-factor order
for the canonical quantity; the
\href{https://github.com/ajt60gaibb/OpenProblemsInNLA/blob/main/references/stepaniants-mf02-2026-09-12/verification/source-review/REVIEW.md}{independent
source review} supplies that synthesis explicitly. No first-discovery or
novelty claim is made for that order or the underlying method.

This self-contained note supplies the explicit comparison in (14), hence
at most one cubic stage per allowed unrestricted multiplication. The
source review did not find that specific comparison in the inspected
publications, but that limited search is not a priority certification.
The result determines the constant-factor asymptotic order of the
canonical MF-02 quantity uniformly in both parameters. The exact
minimum, the optimal leading constant, and the stronger source question
comparing the two errors at the same multiplication budget remain
unanswered. The complete original target and historical references are
retained on the canonical page.

Section 2 proves the Chebyshev comparison and Section 3 proves the cubic
interval estimate. The
\href{https://github.com/ajt60gaibb/OpenProblemsInNLA/blob/main/references/stepaniants-mf02-2026-09-12/verification/independent-review/MF-02-independent-review.md}{independent
mathematical review} reconstructs the complete argument, including the
potentially unattained infima and both endpoint cases. Its
\href{https://github.com/ajt60gaibb/OpenProblemsInNLA/blob/main/references/stepaniants-mf02-2026-09-12/verification/independent-review/exact_algebra_check.py}{exact
checker} confirms 17 rational polynomial and coefficient-sign checks.
These computations supplement the proof. The
\href{https://github.com/ajt60gaibb/OpenProblemsInNLA/blob/main/references/stepaniants-mf02-2026-09-12/README.md}{submission
record} retains both reviews and their exact source fingerprints.
Neither automated-agent review nor these checks constitute external
human peer review or formal verification.

\subsubsection{References}\label{references}

\begin{enumerate}
\def\labelenumi{\arabic{enumi}.}
\tightlist
\item
  N. Amsel et al., \emph{Linear Systems and Eigenvalue Problems: Open
  Questions from a Simons Workshop}, arXiv:2602.05394v3, Section 6.3,
  Problem 6.5 and equations (21)-(22).
  \href{https://arxiv.org/html/2602.05394v3\#S6.SS3}{Primary source}.
\item
  E. H. Rubensson, E. Jarlebring and G. Lorentzon, \emph{Recursive
  expansion of the matrix step function using polynomials of degree
  eight}, arXiv:2606.24701v1, Section 7.
  \href{https://arxiv.org/html/2606.24701v1}{Primary source}.
\item
  J. Chen and E. Chow, \emph{A Stable Scaling of Newton-Schulz for
  Improving the Sign Function Computation of a Hermitian Matrix},
  ANL/MCS-P5059-0114, 2014, Section 3, equations (3.3)-(3.6) and
  Theorems 3.1-3.2.
  \href{https://www.mcs.anl.gov/papers/P5059-0114.pdf}{Primary author
  report}.
\item
  N. Amsel, D. Persson, C. Musco and R. M. Gower, \emph{The Polar
  Express: Optimal Matrix Sign Methods and Their Application to the Muon
  Algorithm}, arXiv:2505.16932v5, 4 May 2026, Theorem 3.1 and Appendix
  F, equation (17). \href{https://arxiv.org/html/2505.16932v5}{Primary
  source}.
\item
  J. H. Cheon, D. Kim and D. Kim, \emph{Efficient Homomorphic Comparison
  Methods with Optimal Complexity}, ASIACRYPT 2020, 221-256, Lemma 3,
  Theorem 3, Table 1 and Section 3.4.
  \href{https://doi.org/10.1007/978-3-030-64834-3_8}{DOI};
  \href{https://eprint.iacr.org/2019/1234.pdf}{author manuscript}.
\end{enumerate}
\end{document}
